\documentclass[12pt, a4paper]{article}

% ── Packages ──────────────────────────────────────────────────────────────────
\usepackage[margin=1in]{geometry}
\usepackage{amsmath, amssymb}
\usepackage{graphicx}
\usepackage{booktabs}
\usepackage{array}
\usepackage{enumitem}
\usepackage{parskip}
\usepackage{titlesec}
\usepackage{hyperref}
\usepackage{cite}

% ── Section formatting ────────────────────────────────────────────────────────
\titleformat{\section}
  {\large\bfseries}{\thesection}{1em}{}

\titleformat{\subsection}
  {\normalsize\bfseries}{\thesubsection}{1em}{}

% ── No paragraph indent, space between paragraphs ─────────────────────────────
\setlength{\parindent}{0pt}
\setlength{\parskip}{6pt}

% ── Hyperlink style ───────────────────────────────────────────────────────────
\hypersetup{
    colorlinks=true,
    linkcolor=black,
    urlcolor=black,
    citecolor=black
}

% ── Document ──────────────────────────────────────────────────────────────────
\begin{document}

% ── Header ────────────────────────────────────────────────────────────────────
\begin{center}
  {\large \textbf{CS4084 Research Project Proposal (RP1)}}\\[10pt]
  {\LARGE \textbf{Performance Analysis of Grover's Search Algorithm\\
  Under Realistic Noise Models on NISQ Simulators}}\\[16pt]

  \textbf{Group Members:}\\[4pt]
  \begin{itemize}[leftmargin=3cm, label=\textbullet]
    \item Saim Haider (22P-9244)
    \item Muhammad Abdullah (22P-9371)
    \item Abdullah (22P-9358)
  \end{itemize}

  \vspace{6pt}
  \textbf{Instructor:} Dr.\ Maqsood M.\ Khan\\
  \textbf{Department:} Computer Science\\
  \textbf{Submission Date:} March 8, 2026
\end{center}

\hrule
\vspace{12pt}

% ── Abstract ──────────────────────────────────────────────────────────────────
\section*{Abstract}

Grover's search algorithm provides a provable quadratic speedup over classical
unstructured search, making it one of the most relevant quantum algorithms for near
term deployment. However, current quantum hardware operates in the Noisy
Intermediate-Scale Quantum (NISQ) regime, where decoherence and gate errors severely
degrade algorithm performance. This research proposes a simulation-based comparative
study of Grover's algorithm under three realistic noise models: depolarizing noise,
bit-flip errors, and phase-flip errors. Using Python and the Qiskit Aer simulator,
Grover's search circuit will be implemented for search spaces of $n = 2, 3, 4, 5$
qubits. Noise will be injected at parameterized error rates swept across the range
$p \in [0.001, 0.1]$ and success probability will be measured at each configuration.
The primary objective is to identify the critical noise threshold $p^{*}$ at which
Grover's algorithm loses its quantum advantage over classical brute-force search, and
to determine how this threshold scales with the number of qubits. Expected outcomes
include empirical noise threshold curves for each noise type, a crossover analysis
comparing Grover's performance against the classical baseline $\frac{1}{N}$, and a
reproducible Qiskit codebase suitable for further research. This work directly
addresses a gap in the existing literature, where no study has systematically
identified this threshold across multiple noise types using modern quantum simulation
tools.

% ── 1. Introduction ───────────────────────────────────────────────────────────
\section{Introduction}

Quantum computing leverages the principles of superposition, entanglement, and
quantum interference to solve certain computational problems significantly faster
than classical computers. Among the most practically relevant quantum algorithms is
Grover's search algorithm, which achieves $\mathcal{O}(\sqrt{N})$ query complexity
for searching an unstructured database of $N$ entries, compared to the classical
$\mathcal{O}(N)$ average case. This quadratic speedup has wide applications in
optimization, cryptanalysis, and database search, making Grover's algorithm a
standard benchmark for evaluating near-term quantum hardware.

Despite its theoretical promise, quantum algorithms face a fundamental challenge in
practice: noise. Physical qubits are susceptible to environmental decoherence, gate
imperfections, and measurement errors. These noise sources collectively corrupt the
quantum interference patterns that Grover's algorithm depends upon to amplify the
probability of measuring the correct answer. Present-day quantum processors operate
in the Noisy Intermediate-Scale Quantum (NISQ) regime, characterized by gate error
rates on the order of $10^{-3}$ to $10^{-2}$ and coherence times in the microsecond
to millisecond range. Under these conditions, even moderately deep circuits such as
Grover's can suffer severe degradation in success probability.

This research proposes a systematic empirical study of Grover's algorithm under
realistic noise conditions using the Qiskit Aer simulator. By sweeping noise error
rates across a controlled range and measuring success probability for each
configuration, this work aims to identify the precise noise threshold beyond which
Grover's algorithm loses its quantum advantage over classical search. The findings
will provide concrete, reproducible benchmarks directly relevant to the development
and calibration of near-term quantum hardware.

% ── 2. Problem Statement ──────────────────────────────────────────────────────
\section{Problem Statement}

\textbf{What specific problem are you solving?}

This research addresses the question: at what noise error rate $p^{*}$ does
Grover's search algorithm lose its quantum computational advantage over classical
brute-force search, and how does this threshold vary across different noise types
(depolarizing, bit-flip, phase-flip) and qubit counts ($n = 2, 3, 4, 5$)?

\textbf{Why does this problem exist?}

Grover's algorithm relies on precise quantum interference to amplify the target
state's probability amplitude. Any noise that disturbs the phase or amplitude of
qubits partially destroys this interference, reducing success probability. As qubit
count increases, the circuit depth of Grover's implementation grows, accumulating
more noise over more gate operations. This makes it critical to understand not just
that noise degrades performance, but precisely how much noise is tolerable before
the algorithm offers no advantage over classical search.

\textbf{What are the limitations of existing approaches?}

Existing studies have addressed this problem only partially. Salas (2008) derived
an analytical noise threshold using Fortran simulation but studied only depolarizing
noise and did not compare against the classical search advantage boundary. Ijaz and
Faryad (2023) used Qiskit Aer to study multiple noise types but focused on comparing
standard versus QSVT implementations rather than identifying the critical threshold
$p^{*}$. Dowarah et al.\ (2025) derived analytical two-threshold models using
Random Matrix Theory but provided no circuit-level simulation. Srivastava et al.\
(2025) proposed quantum switch mitigation techniques but assumed knowledge of the
noise threshold rather than empirically establishing it. No existing study combines
a modern Qiskit simulation framework, multiple isolated noise types, a systematic
error rate sweep, and explicit identification of the quantum advantage loss
threshold.

% ── 3. Literature Review ──────────────────────────────────────────────────────
\section{Literature Review and Base Papers}

\subsection*{Ijaz and Faryad (2023) -- Base Paper}

Ijaz and Faryad conducted a comprehensive noise analysis of Grover's algorithm
using Qiskit Aer for $n = 3, 4, 5$ qubits with 10{,}000 shots per configuration.
Four noise channels were studied: depolarizing, bit-flip (X), phase-flip (Z), and
bit-phase flip (Y). The key finding was that success probability decays
exponentially with error probability $p$ but only linearly with qubit count $n$.
The bit-phase flip channel caused the most severe degradation across all
configurations. This paper is designated as the base paper for our research as it
uses the same Qiskit Aer framework and directly characterizes the noise behavior we
intend to extend.

\textit{Strength:} Real Qiskit Aer simulation with four noise types; exponential
decay relationship confirmed empirically.

\textit{Limitation:} Does not identify the critical threshold $p^{*}$ where
quantum advantage over classical search is lost; focuses on QSVT comparison rather
than threshold analysis.

\subsection*{Salas (2008)}

Salas performed the earliest systematic numerical study of noise in Grover's
algorithm using Fortran-based Monte Carlo simulation for $n = 2$ to $7$ qubits
under a depolarizing channel model. The study derived an exponential damping law
for success probability and an allowed-error law:
$\varepsilon_{\text{th}}(N) \sim N^{-1.1}$, where $N = 2^{n}$. An absolute
threshold of $\varepsilon > 0.043$ was identified above which the algorithm fails
for any qubit count.

\textit{Strength:} First quantitative noise threshold derived for Grover's
algorithm; foundational analytical results still cited in modern literature.

\textit{Limitation:} Uses an outdated Fortran framework predating Qiskit; studies
only depolarizing noise; does not explicitly compare against classical search
advantage.

\subsection*{Dowarah et al.\ (2025)}

Dowarah, Zhang, Khemani, and Kolodrubetz applied Random Matrix Theory (RMT) to
model the noisy Grover operator as a Floquet unitary and derived two analytical
noise thresholds. The computing transition at
$\delta_{c,\text{comp}} \sim \mathcal{O}(2^{-L/2})$ marks where quantum advantage
is lost, and the ergodicity transition at
$\delta_{c,\text{gap}} \sim \mathcal{O}(L^{-1})$ marks where circuit structure
collapses entirely.

\textit{Strength:} State-of-the-art analytical framework; identifies two distinct
thresholds with exact scaling laws.

\textit{Limitation:} Entirely theoretical; restricted to systematic coherent noise;
no Qiskit simulation or empirical validation provided.

\subsection*{Srivastava et al.\ (2025)}

Srivastava, Pati, Chakrabarty, and Bhattacharya proposed using quantum switches
exploiting indefinite causal order to mitigate depolarizing noise in Grover's
algorithm. Two frameworks were developed, with the second framework (delayed
measurement) providing significantly greater noise resilience by preserving
correlations between the input state and quantum switch register throughout all
iterations.

\textit{Strength:} Novel mitigation approach; demonstrates improved success
probability analytically; directly applicable to NISQ devices.

\textit{Limitation:} Entirely analytical with no Qiskit simulation; studies only
depolarizing noise; does not identify the threshold $p^{*}$ for standard Grover
circuits.

\vspace{8pt}

% ── Comparative Table ─────────────────────────────────────────────────────────
\subsection*{Comparative Analysis}

\vspace{4pt}
\begin{table}[htbp]
\centering
\resizebox{\textwidth}{!}{%
\begin{tabular}{p{3.2cm} p{3.2cm} p{4.5cm} p{4.5cm}}
  \toprule
  \textbf{Paper} & \textbf{Method Used} & \textbf{Strength} & \textbf{Limitation} \\
  \midrule
  Ijaz \& Faryad (2023) [Base]
    & Qiskit Aer simulation
    & 4 noise types; exponential decay confirmed
    & No $p^{*}$ threshold; QSVT focus \\[6pt]

  Salas (2008)
    & Monte Carlo, Fortran
    & First threshold law derived
    & Outdated tools; depolarizing only \\[6pt]

  Dowarah et al.\ (2025)
    & Random Matrix Theory
    & Two-threshold analytical model
    & No simulation; coherent noise only \\[6pt]

  Srivastava et al.\ (2025)
    & Analytical, quantum switch
    & ICO mitigation; improved success probability
    & No Qiskit; depolarizing only; no $p^{*}$ \\
  \bottomrule
\end{tabular}}
\end{table}

% ── 4. Proposed Methodology ───────────────────────────────────────────────────
\section{Proposed Methodology}

\textbf{Proposed Algorithm and System}

Grover's search circuit will be implemented from first principles in Qiskit for
$n = 2, 3, 4, 5$ qubits. Each circuit consists of a Hadamard initialization layer,
an oracle operator $\mathcal{O}_\tau = 2|\tau\rangle\langle\tau| - I$ marking the
target state, and a diffusion operator $D = 2|s\rangle\langle s| - I$ performing
inversion about the mean. The circuit applies $k = \lfloor \frac{\pi}{4}\sqrt{N}
\rfloor$ Grover iterations to maximize the ideal success probability:

\begin{equation}
  P_{\text{success}}(k) = \sin^2\!\left((2k+1)\arcsin\frac{1}{\sqrt{N}}\right)
  \approx 1
\end{equation}

\textbf{Noise Injection}

Three noise models will be injected independently using the Qiskit Aer NoiseModel
API. Depolarizing noise applies a completely mixed state with probability $p$.
Bit-flip noise applies a Pauli $X$ gate with probability $p$. Phase-flip noise
applies a Pauli $Z$ gate with probability $p$. Error rates will be swept across
$p \in [0.001, 0.1]$ in steps of 0.005 and each configuration will be executed
with 1024 shots averaged over 5 independent runs.

\textbf{Threshold Identification}

The critical threshold $p^{*}$ is defined as the error rate at which Grover's
measured success probability crosses the classical random search baseline
$\frac{1}{N} = \frac{1}{2^{n}}$. This crossover point will be identified
numerically for each combination of noise type and qubit count.

\textbf{Innovation Over Base Papers}

This study extends Ijaz and Faryad (2023) by restricting analysis to the standard
Grover implementation and explicitly locating $p^{*}$, which their work does not
provide. It extends Salas (2008) by using modern Qiskit tooling and testing
multiple noise channels. It provides the empirical foundation that Srivastava
et al.\ (2025) presuppose but do not establish.

% ── 5. Tools ──────────────────────────────────────────────────────────────────
\section{Tools and Implementation Plan}

\begin{itemize}
  \item \textbf{Programming Language:} Python 3.10+
  \item \textbf{Quantum Framework:} Qiskit 1.x with Qiskit Aer noise simulator
  \item \textbf{Noise API:} \texttt{qiskit\_aer.noise.NoiseModel}
  \item \textbf{Visualization:} matplotlib, numpy
  \item \textbf{Platform:} IBM Quantum Lab (free tier) or local Aer simulator
  \item \textbf{No external dataset required:} The search space is synthetically
        generated as an $n$-qubit register
\end{itemize}

% ── 6. Expected Results ───────────────────────────────────────────────────────
\section{Expected Results}

\textbf{What results do you expect?}

We expect success probability to decay monotonically as error rate $p$ increases
for all noise types. Based on the literature, phase-flip noise is expected to cause
more severe degradation than bit-flip noise for Grover circuits due to the
algorithm's sensitivity to phase interference. Depolarizing noise is expected to
cause the earliest threshold crossing as it simultaneously corrupts both phase and
amplitude information.

\textbf{How will performance be measured?}

Performance will be measured as the probability of measuring the correct target
state across 1024 shots per configuration. The primary output will be a family of
success probability versus error rate curves, one per noise type per qubit count,
with the threshold $p^{*}$ marked on each curve. Secondary metrics will include
circuit depth and gate count per qubit configuration.

\textbf{What improvement do you anticipate compared to base papers?}

Compared to Ijaz and Faryad (2023), this work will produce an explicit quantum
advantage threshold map that their study does not provide. Compared to Salas (2008),
results will be validated on a modern Qiskit circuit simulator with three noise
types rather than one. The threshold values obtained will also serve as a direct
empirical input for noise mitigation strategies such as the quantum switch framework
of Srivastava et al.\ (2025).

% ── 7. Timeline ───────────────────────────────────────────────────────────────
\section{Proposed Timeline}

\begin{itemize}
  \item \textbf{Week 1 to 2:} Literature review and circuit design for all qubit
        configurations
  \item \textbf{Week 3 to 5:} Qiskit circuit implementation and noiseless baseline
        validation
  \item \textbf{Week 6 to 7:} Noise model integration, error rate sweep, and data
        collection
  \item \textbf{Week 8:} Threshold analysis, plot generation, and IEEE paper writing
\end{itemize}

% ── References ────────────────────────────────────────────────────────────────
\begin{thebibliography}{9}

\bibitem{grover1996}
L.~K. Grover, ``A fast quantum mechanical algorithm for database search,''
in \textit{Proc.\ 28th Annual ACM Symposium on Theory of Computing (STOC)},
1996, pp.\ 212--219.

\bibitem{ijaz2023}
M.~A. Ijaz and M.~Faryad, ``Noise analysis of Grover and phase estimation
algorithms implemented as quantum singular value transformations for a small
number of noisy qubits,'' \textit{Scientific Reports}, vol.~13, p.~20144, 2023.

\bibitem{salas2008}
P.~J. Salas, ``Noise effect on Grover algorithm,''
\textit{The European Physical Journal D}, vol.~46, pp.~365--381, 2008.

\bibitem{dowarah2025}
S.~Dowarah, C.~Zhang, V.~Khemani, and M.~Kolodrubetz,
``Phases and phase transition in Grover's algorithm with systematic noise,''
\textit{arXiv preprint arXiv:2406.10344v2}, 2025.

\bibitem{srivastava2025}
S.~Srivastava, A.~K. Pati, I.~Chakrabarty, and S.~Bhattacharya,
``Using quantum switches to mitigate noise in Grover's search algorithm,''
\textit{arXiv preprint arXiv:2401.05866v2}, 2025.

\end{thebibliography}

\end{document}